\documentclass[11pt]{amsart}

\usepackage[T1]{fontenc}
\usepackage{lmodern}
\usepackage{microtype}
\usepackage{mathtools,amssymb,amsthm}
\usepackage[hidelinks]{hyperref}

\hypersetup{
  pdftitle={A four-vertex counterexample to the Clow--van Bommel eternal
            distance-k domination conjecture}
}

\newtheorem{theorem}{Theorem}
\newtheorem{proposition}[theorem]{Proposition}

\begin{document}

\title[A four-vertex counterexample]
{A four-vertex counterexample to the Clow--van Bommel\\
 eternal distance-$k$ domination conjecture}

\date{}

\subjclass[2020]{05C69, 05C57, 05C05}
\keywords{eternal domination, distance-$k$ domination, tree, path,
counterexample}

\begin{abstract}
Let $\gamma_k(G)$ be the distance-$k$ domination number of $G$ and let
$\gamma_{\mathrm{all},k}^{\infty}(G)$ be its eternal distance-$k$ domination
number in the all-guards-move model. Conjecture~8.2 of Clow and van Bommel
asserts that for all $k>2$ and all trees $T$, if
$\gamma_k(T)=\gamma_{\mathrm{all},k}^{\infty}(T)$ then
$\gamma_{\mathrm{all},k}^{\infty}(T)=\gamma_{\lfloor k/2\rfloor}(T)$.
The path $P_4$ at $k=3$ refutes it: there
$\gamma_3=\gamma_{\mathrm{all},3}^{\infty}=1$ while
$\gamma_{\lfloor 3/2\rfloor}=\gamma=2$. The mechanism is parity, so the
statement is false for every odd $k\ge3$, witnessed by every tree of diameter
exactly $k$. The value $\gamma_{\mathrm{all},k}^{\infty}(T)=1$ for
$\operatorname{diam}(T)\le k$ is Lemma 3.3 of the same paper, so no new
mathematics is claimed. Even $k$ is not addressed.
\end{abstract}

\maketitle

\section{The conjecture}

All graphs are finite and simple, and $d(u,v)$ denotes graph distance. For
$k\ge1$ a set $S\subseteq V(G)$ is \emph{distance-$k$ dominating} if every
vertex of $G$ lies within distance $k$ of a vertex of $S$; the least size of
such a set is $\gamma_k(G)$, and $\gamma_1(G)=\gamma(G)$ is the ordinary
domination number. Write $N_k[v]=\{u: d(u,v)\le k\}$.

Following Cox, Meger and Messinger \cite{CMM}, a family $\mathcal F$ of
$g$-element multisets of $V(G)$ (\emph{configurations}) is an \emph{eternal
distance-$k$ dominating family} in the all-guards-move model if
\begin{enumerate}
\item[(i)] every $A\in\mathcal F$ is distance-$k$ dominating, and
\item[(ii)] for every $A\in\mathcal F$ and every $v\in V(G)$ there are
  $B\in\mathcal F$ with $v\in B$ and a bijection of multisets
  $\sigma\colon A\to B$ with $d(a,\sigma(a))\le k$ for all $a\in A$.
\end{enumerate}
The least such $g$ is $\gamma_{\mathrm{all},k}^{\infty}(G)$. Condition~(i) ---
that \emph{every} configuration of the family, not merely the initial one, is
itself distance-$k$ dominating --- is the published definition \cite[\S2]{CMM},
and it is the reading used throughout. The refutation does not depend on this
choice: for the families exhibited below, every configuration of the family is
reachable, so the weaker reading --- requiring only reachable configurations to
distance-$k$ dominate --- gives the same values. Cox, Meger and Messinger also prove the
sandwich
\begin{equation}\label{eq:sandwich}
 \gamma_k(G)\;\le\;\gamma_{\mathrm{all},k}^{\infty}(G)\;\le\;
 \gamma_{\lfloor k/2\rfloor}(G).
\end{equation}

Clow and van Bommel \cite{CvB} prove for $k=2$ that the two outer quantities in
\eqref{eq:sandwich} coincide whenever the two left ones do, and in the
Conclusion of that paper conjecture that the same holds beyond $k=2$:

\begin{quote}
\textbf{Conjecture 8.2} \cite{CvB}. \emph{For all $k>2$ and trees $T$, if
$\gamma_k(T)=\gamma_{\mathrm{all},k}^{\infty}(T)$, then
$\gamma_{\mathrm{all},k}^{\infty}(T)=\gamma_{\lfloor k/2\rfloor}(T)$.}
\end{quote}

\noindent The statement is quoted from the e-print
\texttt{arXiv:2308.00054v1}. That the quotation is faithful, and that no
standing hypothesis of \cite{CvB} has been dropped in transcribing it, is a
documentary matter checked by reading the source, not by the program of \S3.

\section{The counterexample}

\begin{theorem}\label{thm:main}
Let $k=3$ and $T=P_4$, the path $1-2-3-4$. Then
\[
 \gamma_3(P_4)=\gamma_{\mathrm{all},3}^{\infty}(P_4)=1,
 \qquad
 \gamma_{\lfloor 3/2\rfloor}(P_4)=\gamma(P_4)=2 .
\]
Hence the hypothesis of Conjecture~8.2 holds and its conclusion fails, so
Conjecture~8.2 is false.
\end{theorem}

\begin{proof}
$\operatorname{diam}(P_4)=3=k$, so every vertex of $P_4$ is within distance $3$
of every other; each of the four singletons is therefore a distance-$3$
dominating set and $\gamma_3(P_4)=1$. The same computation makes the family
$\mathcal F=\bigl\{\{1\},\{2\},\{3\},\{4\}\bigr\}$ of all four singletons
eternal: each member distance-$3$ dominates, and an attack at $v$ is answered by
moving the lone guard to $v$, a move of length at most $3$. So
$\gamma_{\mathrm{all},3}^{\infty}(P_4)=1$, the value $0$ being impossible.
Finally $\lfloor 3/2\rfloor=1$, and no single vertex of $P_4$ dominates it:
$N_1[1]=\{1,2\}$, $N_1[2]=\{1,2,3\}$, $N_1[3]=\{2,3,4\}$, $N_1[4]=\{3,4\}$.
Since $\{2,3\}$ dominates, $\gamma(P_4)=2$. The conclusion of Conjecture~8.2
would require $1=2$.
\end{proof}

Write $\operatorname{rad}(T)$ for the radius; in a tree
$\operatorname{rad}(T)=\lceil\operatorname{diam}(T)/2\rceil$.

\begin{proposition}\label{prop:oddk}
Let $k\ge3$ be odd and let $T$ be any tree with $\operatorname{diam}(T)=k$.
Then $\gamma_k(T)=\gamma_{\mathrm{all},k}^{\infty}(T)=1$ while
$\gamma_{\lfloor k/2\rfloor}(T)\ge2$, so $T$ is a counterexample to
Conjecture~8.2 at that $k$.
\end{proposition}

\begin{proof}
If $\operatorname{diam}(T)\le k$ then a single vertex distance-$k$ dominates
$T$ and every move is legal, so exactly as in Theorem~\ref{thm:main}
$\gamma_k(T)=\gamma_{\mathrm{all},k}^{\infty}(T)=1$; the eternal value is also
the second case of Lemma~3.3 of \cite{CvB}, proved there five sections before
Conjecture~8.2. Next $\gamma_j(T)=1$ if and only if $\operatorname{rad}(T)\le
j$. With $\operatorname{diam}(T)=k$ odd, $\operatorname{rad}(T)=(k+1)/2>
(k-1)/2=\lfloor k/2\rfloor$, so $\gamma_{\lfloor k/2\rfloor}(T)\ge2$ and the
conclusion fails.
\end{proof}

The obstruction is parity: $2\lfloor k/2\rfloor+1$ equals $k+1$ for even $k$ but
only $k$ for odd $k$. The least tree of diameter exactly $k$ is $P_{k+1}$, which
gives Theorem~\ref{thm:main} at $k=3$.

\medskip
\noindent\textbf{Scope.} Conjecture~8.2 is universally quantified over $k>2$, so
a single odd $k$ refutes it as written. The mechanism needs $k$ odd, and nothing
above bears on the restriction of Conjecture~8.2 to even $k$, or on the other
conjectures and questions of \cite[\S8]{CvB}. We refute the printed statement
and do not speculate about intent.

\medskip
\noindent\textbf{Attribution.} Theorem~\ref{thm:main} is Lemma~3.3 of
\cite{CvB} --- quoted, not reproved --- together with $\gamma(P_4)=2$; the
observation that these contradict Conjecture~8.2, with the parity account of
it, is what this note adds.

\section{Verification}

Everything above is checkable by hand from the objects printed here: $P_4$, and
$P_{k+1}$ for Proposition~\ref{prop:oddk}. The accompanying program
\texttt{verify.py}, with its recorded transcript \texttt{verify.output.txt},
re-derives the displayed numbers independently --- distances, radii and
diameters by breadth-first search, each $\gamma_j$ by an exact minimum set
cover, and each $\gamma_{\mathrm{all},k}^{\infty}$ by a greatest-fixed-point
solution of the game of \S1 over its full configuration space. It uses exact
integer arithmetic and the Python standard library only.

The transcript is that of an earlier and wider draft, and the reader should
read it accordingly. Its Steps 1, 2 and 4 are the checks that belong to this
note: Step 1 and Step 2 verify every number in Theorem~\ref{thm:main}, and
Step 4 instantiates Proposition~\ref{prop:oddk} at $k=3,5,7$ and on one
non-path tree, the general statement resting on the proof above and not on
those instances. Everything else in the transcript concerns material this note
does not contain and does not claim: Step 3 treats $P_8$ at $k=3$; Steps 5
and 6 are headed ``Corollary 4'' and ``Corollary 5'' and refer to a path family,
to published closed forms and to a correction to \cite{CMM}, none of which is
stated here; Step 7 is a bounded sweep of the $95$ trees of order at most $9$;
and the closing scope notes mention a wider census of trees, and locations
inside the source tarball, that this note does not cite. The internal numbering
of the transcript is that of the earlier draft and does not match the numbering
printed above. The program checks no statement about the contents of \cite{CvB}
or \cite{CMM}; those are documentary claims, checked by reading.

\begin{thebibliography}{9}

\bibitem{CvB}
A.~Clow and C.~M.\ van Bommel,
\emph{Eternal distance-2 domination in trees},
Ars Combinatoria \textbf{163} (2025), 29--49.
\href{https://doi.org/10.61091/ars163-03}{doi:10.61091/ars163-03};
\href{https://arxiv.org/abs/2308.00054v1}{arXiv:2308.00054v1}.
Conjecture~8.2 and Lemma~3.3 are quoted from the e-print.

\bibitem{CMM}
D.~Cox, E.~Meger, and M.~E.\ Messinger,
\emph{Eternal distance-$k$ domination on graphs},
La Matematica \textbf{2} (2023), no.~2, 283--302.
\href{https://doi.org/10.1007/s44007-023-00044-3}{doi:10.1007/s44007-023-00044-3};
\href{https://arxiv.org/abs/2104.03835}{arXiv:2104.03835}.
Source of the definition in \S1 and of the sandwich \eqref{eq:sandwich}.

\end{thebibliography}

\end{document}
